% ================================================================
%  CHAPTER 6 -- DISCUSSION AND CONCLUSION
%  Status: DRAFT v1 (2026-05-18).
% ================================================================
\chapter{Discussion and Conclusion}
\label{ch:conclusion}

% ----------------------------------------------------------------


This thesis began from a single observation: training a neural Bayes
estimator involves Monte Carlo integration of the Bayes risk, and the
gradient that drives training is therefore a Monte Carlo estimator
with a variance. That variance is a property of how the loss is
sampled, not of the inference problem, and it can be reduced.

The primary reduction attempted here is the Rao-Blackwellised loss. When a
block of the parameter vector has a tractable conditional posterior,
that block is integrated out of the loss analytically rather than
sampled. For squared-error loss this amounts to substituting the
conditional posterior mean for the sampled target, with the residual
conditional variance entering only as a constant that does not affect
the gradient. The resulting training gradient is the
Rao-Blackwellisation of the standard one, and the law of total
variance gives a positive semi-definite gradient-variance gap: the
technique can never increase the variance of the training gradient.
For a linear network the guarantee strengthens to the fitted weights
themselves.

The empirical study on Bayesian linear regression bore this out. The
Rao-Blackwellised estimator attained a lower loss than the standard
estimator in every cell of the experimental grid, closed about half
the gap between the standard estimator and the Bayes-optimal floor at
the baseline, and matched the standard estimator's accuracy at a
quarter of the batch size.


Chapter~\ref{ch:importance-sampling} placed Rao-Blackwellisation
alongside importance sampling, a second variance-reduction technique
acting on the training of the estimator. The comparison is the
clearest summary of what makes Rao-Blackwellisation attractive.
Importance sampling can in principle sharpen estimation by
oversampling the regions that contribute most to the loss, but only
when the bottleneck on training is gradient noise and the proposal is
matched to it. On the autoregressive model neither condition held: the
trained estimator already sat close to the Bayes-optimal floor across
the prior, and the proposal oversampled the very regions it handled
best, so reweighting found no error worth recovering and, pushed hard,
only added gradient noise.
Rao-Blackwellisation carries an unconditional guarantee, the
variance gap being positive semi-definite for every model and
architecture, at the cost of requiring a tractable conditional
posterior. Where that conditional structure is present, as it is in
any conditionally conjugate hierarchical model, Rao-Blackwellisation
is the safer and stronger of the two.

One distinction underlies this comparison and is worth stating
directly. Reducing the variance of the loss --- or, more precisely, of
the gradient that drives training --- is not really the primary goal;
the goal is reducing the variance of the estimate that training
produces, the fitted weights $\hat{\bm{\gamma}}$ and the estimator they
define. A cleaner gradient is valuable only insofar as it yields a less
variable estimator. At this level the two techniques separate most
sharply. For a linear network, Section~\ref{sec:linear-network} showed
that Rao-Blackwellisation reduces not only the variance of the gradient
but the variance of the fitted weights themselves, with the last-layer
argument suggesting the same holds approximately for nonlinear
networks. Importance sampling admits no such statement: it bears, and
only contingently, on the variance with which a minibatch estimates the
gradient, and nothing in it carries through to the variance of the
estimator that training returns. That a guarantee can be carried to the
fitted weights for one technique and not the other is the deepest form
of the contrast between them.


Several limitations bound the results. The Rao-Blackwellised loss
requires a tractable conditional posterior for the integrated block;
where no parameter block is conditionally conjugate the construction
does not apply, and only the sampled block benefits when one block is.
The experimental evidence is confined to two models (linear
regression with a multilayer perceptron, and an autoregressive process
with a recurrent network) and does not establish how the method
behaves for richer simulators or larger architectures. The estimators
studied are point estimators and provide no quantification of
posterior uncertainty.

% ----------------------------------------------------------------
\section{Future work}
\label{sec:future-work}

Several directions follow naturally. The linear-network result of
Section~\ref{sec:linear-network} suggests, but does not prove, a
corresponding statement for fully nonlinear networks; making the
last-layer argument rigorous, or bounding its error, would close that
gap. The squared-error construction of Section~\ref{sec:rb-mse} could
be extended to other losses (the quantile loss of
Section~\ref{sec:loss-summaries}, for instance) to Rao-Blackwellise
estimators of posterior quantiles rather than means. Where a parameter
block is only approximately conjugate, a partial or approximate
Rao-Blackwellisation may still reduce variance, and quantifying the
resulting bias--variance trade-off is an open question. The
variance-reduction toolbox of Monte Carlo contains further techniques
not explored here (control variates, antithetic sampling, and
stratified proposals among them), and a systematic comparison of
these against Rao-Blackwellisation for estimator training, including
how they might be combined, is a natural direction. A related
direction is the reparameterisation of the estimator's outputs to
decouple weakly identified parameters, an alternative to variance
reduction for the boundary difficulties of
Chapter~\ref{ch:importance-sampling}. Most importantly, the method
should be tested on the larger, less tractable simulators that are the
natural setting for simulation-based inference.

% ----------------------------------------------------------------
\section{Conclusion}
\label{sec:final}

Neural Bayes estimators amortise Bayesian inference by replacing
per-dataset sampling with a single trained mapping. Because their
training is itself Monte Carlo integration, the variance-reduction
toolbox of Monte Carlo applies to the training loop. This thesis has
shown that Rao-Blackwellisation, in particular, is a cheap and
reliable lever: where a parameter block is conditionally conjugate,
integrating it out of the loss provably reduces the variance of the
training gradient, and in the regimes studied it drives the trained
estimator substantially closer to the Bayes-optimal floor. Conditional
conjugacy, long exploited to make Gibbs sampling possible, turns out
to be equally useful for training the estimators that aim to replace
it.
